\section{Eucharistic Discipline: Reception, Reverence, and Ministry}

The Eucharist is not made safer by collapsing law into taste. Communion on the tongue is a universal right under current Roman discipline; Communion in the hand is a lawful territorial permission where recognized and may be chosen by the communicant. Kneeling expresses adoration; standing can lawfully express the ordered procession of the pilgrim Church. An extraordinary minister can distribute validly consecrated Communion when the Church deputes that service; routine use without genuine need contradicts the discipline that makes the ministry extraordinary. None of those propositions cancels another.

\subsection{Who should receive}

Canon 912 begins generously: any baptized person not prohibited by law can and must be admitted to Holy Communion. The surrounding canons give the sacramental shape of that generosity. A person conscious of grave sin must not receive without prior sacramental confession unless a grave reason and lack of opportunity concur, in which case perfect contrition includes the resolve to confess as soon as possible (can.~916). Canon 915 concerns the minister's public discipline toward those under a canonical exclusion; it is not a license for ministers to investigate consciences in the Communion line.

The ordinary fast is at least one hour before Communion from food and drink, except water and medicine; elderly, infirm, and caregivers receive the law's mitigation (can.~919). A person may receive twice on the same day, but the second reception must occur within a Mass in which the person participates, apart from the special law of Viaticum (can.~917 and its authentic interpretation). Children require sufficient knowledge and careful preparation, while danger of death requires only the capacity to distinguish Christ's Body and receive reverently (can.~913).

These are minimum juridical boundaries, not a complete spirituality. Recollection, appropriate dress, sacramental confession, thanksgiving, restitution, and reconciliation with one's neighbor dispose the person to receive fruitfully. Neither Missal teaches automatic Communion. The older order's repeated \latin{Domine, non sum dignus} and the newer order's common form both place mercy before the communicant.

\subsection{Tongue and hand: history and present law}

Patristic sources witness reception in the hand under careful disciplines, including Cyril of Jerusalem's mystagogical instruction. Other early evidence stresses fear of losing a particle, reception by the mouth, and changing customs. The later Western norm of the minister placing the Host on the tongue developed with intensified Eucharistic reverence, kneeling, and protection of the Species. It is bad history to call hand reception a Protestant invention without ancient precedent; it is equally bad history to import the modern standing queue, bare-handed distribution, and surrounding culture wholesale into one patristic passage.

Paul VI's \emph{Memoriale Domini} (1969) surveyed the Latin episcopate, retained the traditional manner of ministerial placement on the tongue, and warned of dangers to reverence, doctrine, and fragments. It nevertheless allowed a controlled territorial exception where the contrary practice had already taken root and where the required episcopal vote and Apostolic confirmation were obtained. A 1974 notification elaborated the manner of the permission. The United States later received it. Communion in the hand is therefore neither the universal default created by Vatican II nor an intrinsically sacrilegious act when lawfully and reverently chosen.

Current universal discipline, as summarized by \emph{Redemptionis Sacramentum} 90--92, forbids denial solely because a communicant kneels; every Catholic has the right to receive on the tongue; and the communicant may receive in the hand where the conference has the recognized permission. If there is a danger of profanation, Communion in the hand is not to be given. The current United States GIRM ordinarily names standing as the common posture and a bow as the act of reverence, but expressly says a kneeling communicant is not to be denied on that ground. The communicant, not the distributor, chooses tongue or hand where both are lawful.

\begin{fourstable}{Claim}{Current judgment}{Reason}{Error avoided}
“Tongue is forbidden because the parish stands.” & False. Reception on the tongue remains a right, and kneeling alone is not grounds for refusal. & Universal norms control territorial posture. & Local uniformity cannot erase a recipient's protected mode.\\
“Hand reception is always sacrilege.” & False as a universal moral judgment where competent authority permits it and the Host is received reverently. & The Church can regulate the mode of reception; sacrilege requires profanation or grave contempt, not mere use of a permitted mode. & Prudential preference is not elevated into a validity or orthodoxy test.\\
“The early Church proves the modern practice should be universal.” & False. & Ancient evidence is diverse and embedded in disciplines unlike the modern practice; antiquity alone is not a norm. & Archaeologism does not replace living authority.\\
“A dropped particle does not matter because Christ is wholly present.” & False conclusion. & Christ is wholly present under every remaining Eucharistic species; that is precisely why fragments receive serious care. & Metaphysical truth strengthens rather than relaxes reverence.\\
\end{fourstable}

In the 1962 Order, the normative ritual is kneeling reception on the tongue from the priest, with the Communion plate used to guard the Host. That celebration must follow its own approved books. A general contemporary permission for hand reception or lay distribution does not by itself authorize a priest to hybridize the 1962 rite; canon 846 forbids private addition, omission, or alteration. A contrary particular faculty would need to be established, not presumed.

\subsection{Ordinary and extraordinary ministers}

Terminology carries theology. The ordinary ministers of Holy Communion are bishop, priest, and deacon (can.~910 \S1). An instituted acolyte and another faithful deputed according to canon 230 \S3 are extraordinary ministers of Holy Communion. The unordained person is not an “extraordinary minister of the Eucharist”: only a validly ordained priest can confect the Eucharistic sacrifice, and only bishop or priest acts as its minister in that strict sense.

\emph{Redemptionis Sacramentum} 151--160 describes extraordinary distribution as supplementary and provisional. It is warranted where no ordained minister is present, where the ordained minister is impeded by weakness, advanced age, or another genuine reason, or where the number of communicants would unduly prolong the celebration. It rejects use merely to increase lay participation, routine substitution while sufficient clergy are present, and an expansive reading of slight prolongation. The diocesan bishop regulates deputation within universal law; a pastor does not create a parallel lay office by scheduling custom alone.

An extraordinary minister distributes the already consecrated Species validly when lawfully deputed. Improper use does not invalidate the Mass, undo consecration, or make a recipient consume bread. It can be a serious liturgical abuse, obscure the relation of Orders to Eucharist, expose the Sacrament to risk, and require correction. Conversely, genuine necessity does not demean the Eucharist: the Church's deputation orders humble service to sacramental access.

\begin{longtable}{>{\bfseries\raggedright\arraybackslash}p{0.2\linewidth}>{\raggedright\arraybackslash}p{0.34\linewidth}>{\raggedright\arraybackslash}p{0.38\linewidth}}
\toprule
\textbf{Circumstance} & \textbf{Disciplinary result} & \textbf{Sacramental result}\\
\midrule
\endhead
Priest and deacon can distribute without unreasonable prolongation & They should exercise their ordinary ministry; adding lay distributors merely for symmetry or speed is not justified. & The Eucharist remains valid if extraordinary ministers are nevertheless used; the issue is liceity and sign.\\
Large congregation and too few non-impeded ordinary ministers & Properly instituted or deputed faithful may assist according to diocesan norms. & Their distribution does not consecrate; it conveys the already consecrated Body and Blood.\\
Priest remains seated while lay ministers distribute without impediment & Contrary to the supplementary nature of the ministry. & Does not by itself invalidate the Mass or Communion.\\
No priest is available for Mass & A lawful Communion service may sometimes be led under diocesan norms using reserved Sacrament. & No Eucharistic sacrifice occurs and no new consecration is effected. It is not “Mass without a priest.”\\
1962 Mass & The proper books assign distribution to the priest and provide no lay extraordinary-minister role. & A modern lay role may not be inserted by private initiative.\\
\bottomrule
\end{longtable}

\subsection{One species, both species, and intinction}

Christ is whole and entire---Body, Blood, soul, and divinity---under either Eucharistic species and under every part after separation. Communion under bread alone is therefore complete and is the inherited discipline of the 1962 parish rite. The postconciliar Missal widens Communion under both kinds because the fuller sign more clearly expresses the Eucharistic banquet, new covenant in Christ's Blood, and eschatological feast. Fuller sign does not mean more Christ or more grace by quantity.

The diocesan bishop and liturgical law govern when the chalice is offered. A communicant is never required to receive from it, and pastoral risk, numbers, health, or danger of profanation can counsel against it. The communicant may not self-communicate by taking the Host or passing the chalice. Intinction is performed by the minister; the communicant receives the intincted Host on the tongue and may not dip a Host by hand. Bread used must remain valid wheat bread, and wine valid grape wine.

For serious gluten intolerance, low-gluten hosts that remain truly wheat bread can be valid; completely gluten-free substitutes cannot. Mustum that remains grape product can be used under competent permission for a priest or communicant unable to take ordinary wine. These provisions join sacramental realism to pastoral care rather than altering matter by sentiment.

\subsection{Reservation, vessels, fragments, and the sacred place}

Consecrated Species are reserved principally for Viaticum and Communion of the sick, and secondarily for adoration. The tabernacle must be secure and worthy; the sanctuary lamp signifies the Presence. Sacred vessels, purification, transport, and custody follow law because Christ remains present as long as the Eucharistic species remain. Neither casual handling nor the opposite superstition that grace depends on a particular vessel understands the reason for discipline.

The current books permit designated ministers to purify vessels under their law; the 1962 books assign their own roles and actions. Spills and dropped Hosts require immediate reverent response according to the applicable rite. The minister should watch that the communicant consumes the Host and should prevent removal or profanation. These are not aesthetic extras. They are practical consequences of transubstantiation.

\subsection{Lawful minimum and ars celebrandi}

Discipline cannot generate devotion by decree, but it can protect the conditions in which devotion grows. Tongue reception should never become a badge of superiority; hand reception should never be handled casually. Kneeling should not become defiance; standing should not become indifference. Clergy should not treat extraordinary ministers as a staffing convenience; the faithful should not treat genuine deputation as counterfeit priesthood.

At both Missals the strongest reform begins before arguments about posture: preach the Real Presence, restore confession, observe the fast, celebrate without haste, use worthy music and silence, protect every fragment, and give thanksgiving time. The controversy becomes fruitful when the object is Christ rather than the identity of a faction.

\governingjudgment{The received-on-the-tongue discipline has unique protective and symbolic strengths and remains every communicant's right. Lawful hand reception is not intrinsically invalid, heretical, or sacrilegious; it demands vigilant reverence and can be withheld where profanation threatens. Extraordinary lay distribution is truly extraordinary: necessity can justify it, routine convenience cannot, and misuse is an abuse rather than a retroactive invalidation of the Eucharist.}
